\section{The Contraction of Tight Zigzags}

A zigzag is contracted if non of its interior morphisms are support idempotents. More specifically, 
a zigzag $F=\zigzagTuple{f}{n}$ is contracted iff
(i) for each $i$, $2 \leq i \leq n-1$, it is not the case that $\rst{f_i}={f_i}$ and (ii) for each $i$, $1 \leq i \leq n-1$, it is not the case that $\rst{f'_i}={f'_i}$.

\subsection{One step contraction}
\subsubsection{Rule for removal of a forward morphism}
If $F=\zigzagTuple{f}{n}$ is a zigzag and the morphism $f_i$ is an interior support idempotent then we define $F$ with $f_i$ removed to be the zigzag
$\tuple{f_1,f'_1,..f'_{i-1},f'_i \circ f_i \circ f'_{i+1},f_{i+1},...f_n}$
as shown in this diagram 

Diagramatically: if $f_i$ is a support idempotent 
(i.e. if $\rst{f_i} = f_i = \rng{f_i}$) then
\begin{wrappedZigzag}{x}{y}{f}
    \SetIndexTranslation{1}{i-1}
    \SetIndexTranslation{2}{i}
    \SetIndexTranslation{3}{i+1}
    \arraycolsep=12pt
    \zigzagThree{x}{y}{f}
\end{wrappedZigzag}
contracts to
\begin{wrappedZigzag}{x}{y}{f}
    \SetIndexTranslation{1}{i-1}
    \SetIndexTranslation{2}{i+1}
    \SetObjectContent{x2}{\kern0.4cm x_{i+1}\kern0.3cm} 
    \SetMorphismContent{fprime1}{f'_{i-1}\circ f_i \circ f'_i}
    \SetMorphismLabelPosition{f2}{b} %flips the label out of the way 
                                     % of the lengthy fprime1 label
    \arraycolsep=10pt
    \zigzagTwo{x}{y}{f}
\end{wrappedZigzag}
which in the case that $F$ is tight simplifies to 
\begin{wrappedZigzag}{x}{y}{f}
    \SetIndexTranslation{1}{i-1}
    \SetIndexTranslation{2}{i+1}
    \SetObjectContent{x2}{\kern0.2cm x_{i+1}\kern0.1cm} 
    \SetMorphismContent{fprime1}{f'_{i-1} \circ f'_i}
    \SetMorphismLabelPosition{f2}{b} %flips the label out of the way 
                                     % of the lengthy fprime1 label
    \arraycolsep=10pt
    \zigzagTwo{x}{y}{f}
\end{wrappedZigzag}


If $F$ is tight then it follows that 
$F$ with $f_i$ removed is    $\tuple{f_1,f'_1,..f'_{i-1},f'_i \circ f'_{i+1},f_{i+1},...f_n}$.

\subsubsection{Rule for removal of a reverse morphism}

If $F=\zigzagTuple{f}{n}$ is a zigzag 
and the morphism $f'_i$ is a interior support idempotent 
then we define $F$ with $f'_i$ removed to be the zigzag
$\tuple{f_1,f'_1,..f'_{i-1},f_i \circ f'_i \circ f_{i+1},f'_{i+1},...f_n}$. 
If $F$ is tight then it follows that 
$F$ with $f_i$ removed is    $\tuple{f_1,f'_1,..f'_{i-1},f_i \circ f_{i+1},f'_{i+1},...f_n}$.

Diagramatically: if $f'_i$ is a support idempotent  then
\begin{wrappedZigzag}{x}{y}{f}
    \SetIndexTranslation{1}{i-1}
    \SetIndexTranslation{2}{i}
    \SetIndexTranslation{3}{i+1}
    \SetIndexTranslation{4}{i+2}
    \arraycolsep=12pt
    \zagzigThree{x}{y}{f}
\end{wrappedZigzag}
contracts to
\begin{wrappedZigzag}{x}{y}{f}
    \SetIndexTranslation{1}{i-1}
    \SetIndexTranslation{2}{i+1}
    \SetIndexTranslation{3}{i+2}
    \SetObjectContent{x2}{\kern8pt x_{i+1}\kern3pt} 
    \SetMorphismContent{f2}{f_i \circ f'_i \circ f_{i+1}}
    \SetMorphismLabelPosition{fprime1}{a} %flips the label out of the way 
                                     % of the lengthy f1 label
    \arraycolsep=12pt
    \zagzigTwo{x}{y}{f}
\end{wrappedZigzag}
which in the case  that $F$ is tight simplfies to
\begin{wrappedZigzag}{x}{y}{f}
    \SetIndexTranslation{1}{i-1}
    \SetIndexTranslation{2}{i+1}
    \SetIndexTranslation{3}{i+2}
    \SetObjectContent{x2}{\kern2pt x_{i+1}\kern1pt} 
    \SetMorphismContent{f2}{f_i \circ  f_{i+1}}
    \SetMorphismLabelPosition{fprime1}{a} %flips the label out of the way 
                                     % of the lengthy f1 label
    \arraycolsep=12pt
    \zagzigTwo{x}{y}{f}
\end{wrappedZigzag}

\subsection {Confluence}
\newcommand{\contract}[1]{contract(#1)}

If $F$ is a tight zigzag then, by the rules described, 
it can be successively contracted by the elimination of any interior morphism which is a support idempotents. Eventually
a zigzag having no interior support morphisms is reached. 

To establish that the contraction is unique, by Newman's lemma, 
it is suffiecient 
to establish the following lemma:
\begin{lemma}
If Z is a zigzag to which two different contractions apply yielding
$Z_1$ and $Z_2$ as immediate contractions of $Z$ then there exists
a zigzag $Z'$ to which both $Z_1$ and $Z_2$ can be contracted through a series of zero, one or more contractions. 
\end{lemma}
\begin{proof}
This is clearly the case if the two morphisms 
that can be eliminated are apart in $Z$ for then the eliminations do not interfere and applying in them in either order we arrive at the same zizgzag. 
If the support idempootents to be eliminated are not apart in $Z$ 
--- if they are two successive morphisms of Z --- 
then the elimination of either 
removes the possibility of eliminating the other. 

Suppose successive morphisms $z_i$ and $z'_i$ are support idempotents.

\begin{wrappedZigzag}{x}{y}{z}
\SetIndexTranslation{1}{i-1}
\SetIndexTranslation{2}{i}
\SetIndexTranslation{3}{i+1}
\arraycolsep=12pt
\zigzagThreePointFive{x}{y}{z}
\end{wrappedZigzag}
Contracting by elimination of $z_i$ we get
\begin{wrappedZigzag}{x}{y}{z}
\arraycolsep=12pt
\SetIndexTranslation{1}{i-1}
\SetIndexTranslation{2}{i}
\SetIndexTranslation{3}{i+1}
\SetMorphismContent{zprime1}{z'_i \circ  z'_{i-1}}
                                 % make room for lengthy f'1 label
\SetMorphismLabelPosition{z2}{b} %flip the next label out of the way 
\SetObjectContent{x2}{\kern2pt x_{i+1}\kern1pt} % add a bit of space
\zigzagTwoPointFive{x}{y}{z}
\end{wrappedZigzag}
whereas contracting by elimination of $z'_i$ we get
\begin{wrappedZigzag}{x}{y}{z}
\arraycolsep=12pt
\SetIndexTranslation{1}{i-1}
\SetIndexTranslation{2}{i}
\SetIndexTranslation{3}{i+1}
\SetMorphismContent{z2}{z_i \circ  z_{i+1}}
                                 % make room for lengthy f2 label
\SetMorphismLabelPosition{zprime1}{a} %flip the previous label 
                                     % of the lengthy f'1 label
\SetObjectContent{x2}{\kern2pt x_{i+1}\kern1pt} % add a bit of space
\zigzagTwoPointFive{x}{y}{z}
\end{wrappedZigzag}
but these are identical zigzags 
because from the idempotency of $z_i$ and $z'_i$ 
and from the tightness of $Z$ it follows that
\[
z'_{i-1} \circ z'_i = z'_{i-1}
\]
and
\[
z_i \circ z_{i+1} = z_{i+1}.
\]
So in this case no further eliminations  required 
to reach a common $Z'$ --- we can take  $Z'= Z_1=Z_2$.

The remaining case where the successive morphisms are $z'_{i-1}$ and $z_i$ is taken account of in the same way and completes the proof that contraction by elimination of support idempotents is confluent. 
\end{proof} 

Following this it is possible to speak of 
the contraction of a tight zigzag
and for the contraction of the tight zigzag $Z$ we will write $\contract{Z}$.

\subsection{Composition of Contracted Tight Zigzags}

From the description of contraction we see the following
\begin{observation}
\label{observationRegardingTheLeadingMorphimOfContractedTightZigzag}
For any tight zigzag $Z=\zigzagTuple{z}{n}$, the leading morphism
of the contrcated zigzag $\contract{Z}$ is $z_1 \circ z'$, for some morphism $z'$. 
\end{observation}

\begin{lemma}
\llabel{contractionPushSnug}
For any tight zigzags $F$ and $G$ such that $F \cnc G$ is defined 
and that fit together snugly,
\[
\contract{F\cnc G} = \contract{F\cnc \contract{G}} 
\]
\end{lemma}
\begin{proof}
Suppose $F=\zigzagTuple{f}{n}$ and that $G=\zigzagTuple{g}{m}$.
 By the observation above the first morphism of $\contract{G}$ is 
$g_1 \circ g'$ for some $g'$.  
But the n'th forward morphism of $F \cnc G$ is $f_n \circ g_1$ and forward and reverse morphsims  after this position are exactly as those of $G$ that succeed $g_1$. Therefore the same eliminations that are followed to reach $\contract{G}$ from $G$ apply to  $F \cnc G$. 
In the resulting intermediate contraction of $F \cnc G$
 the n'th forward morphism is 
$f_n \circ g \circ g'$, the morphims that preceded this
are the non-final morphisms of $F$, the morphisms that succeed
it are the non-initial morphisms of $\contract{G}$.
The intermediate contraction of $F \cnc G$ is therefore  
 $F \cnc \contract{G}$. Since this is an intermediate contraction, contracting further, we get 
 $\contract{F\cnc \contract{G}} =  \contract{F\cnc G}$, as required.        
\end{proof}

\begin{lemma}
\llabel{newlemmas}
(i) If $F$ is a tight zigzag then 
\[
\rst{\contract{F}}= \rst{F}
\]
\[
\rng{\contract{F}}= \rng{F}
\]
(ii) If $Z$ is any zigzag then \highlight{is the rhs defined?}
\[
\contract{\tighten{Z}} = \contract{\tighten{\contract{Z}}}
\]
(iii) If $G$ is a tight zigzag then \highlight{???}
\[
\contract{\rng{F} \circ G} = \rng{F} \circ \contract{G}
\]
\end{lemma}
\begin{proof}
\end{proof}


\begin{worktt}
\highlight{Transit} from snug and $\cnc$ to $\circ$.
shift as early as possible. 
\end{worktt}
\begin{lemma}
\llabel{contractionPushGeneric}
For any tight zigzags $F$ and $G$ such that $F \cnc G$ is defined 
\oldt{and that fit together snugly},
\[
\contract{\tighten{F\cnc G}}
 = \contract{\tighten{F\cnc \contract{G}}} 
\]
\end{lemma}
\begin{proof}
We need various new lemmas
\begin{align*}
\contract{\tighten{F\cnc G}}
 &= \contract{
                \tighten{F \cnc \rst{G}}
                \cnc
                \tighten{\rng{F} \cnc G}
             }
             &&\mbox{existing lemma,}\\
 &= \contract{
                \tighten{F \cnc \rst{G}}
                \cnc
                \contract{\tighten{\rng{F} \cnc G}}
             }
             &&\mbox{existing lemma,}\\
 &= \contract{
                \tighten{F \cnc \rst{G}}
                \cnc
                \contract{\tighten{\contract{\rng{F} \cnc G}}}
             }
             &&\mbox{new lemma,}\\
 &= \contract{
                \tighten{F \cnc \rst{G}}
                \cnc
                \tighten{\contract{\rng{F} \cnc G}}
             }
             &&\mbox{existing lemma,}\\
 &= \contract{
                \tighten{F \cnc \rst{G}}
                \cnc
                \tighten{\rng{F} \cnc \contract{G}}
             }
             &&\mbox{a different new lemma,}\\
 &= \contract{
                \tighten{F \cnc \rst{\contract{G}}}
                \cnc
                \tighten{\rng{F} \cnc \contract{G}}
             }
             &&\mbox{yet another new lemma,}\\
 &= \contract{\tighten{F\cnc \contract{G}}} 
             &&\mbox{existing lemma, as required.}\\
\end{align*}
\end{proof}

From the earlier lemma \lref{contractionPushSnug} we have that
\begin{lemma}
\llabel{snugAssociativity}
If tight zigzags $F$, $G$ and $H$ fit together snugly then
\[
\contract{\contract{F \cnc G} \cnc H}
= \contract{F \cnc \contract{G \cnc H}}
\]
\end{lemma}
\begin{proof}
By applying the previous lemma twice:
\begin{align*}
\contract{\contract{F \cnc G} \cnc H}
&= \contract{F \cnc G \cnc H} \\
&= \contract{F \cnc \contract{G \cnc H}}
\end{align*}
\end{proof}
following which we can drop the assumption the $F$, $G$ and $H$ fit together snugly:

\begin{worktt}
\begin{lemma}
\llabel{associativityOfContractedComposition}
For any tight zigzags $F$, $G$ and $H$ such that $F \circ G \circ H$ is defined
\[
\contract{\tighten{
                   \contract{\tighten{F \cnc G}} \cnc H
                  }
         }
= \contract{\tighten{
                      F \cnc \contract{\tighten{G \cnc H}}
                     }
            }
\]
\end{lemma}
\begin{proof}
Use lemma {howtotightenaThreefoldconcatenation}?
\end{proof}
\end{worktt}
We will define composition in the category of contracted tight zigzags 
in terms of composition in the category of tight zigzags.
If $F$ and $G$ are tight zigzags then we define the composition 
as 
\newcommand{\circC}{\mathrel{\underset{c}{\circ}}}
\newcommand{\circT}{\mathrel{\underset{ct}{\circ}}}
\[
F \circC G = \contract{F \circT G} = \contract{\tighten{F \cnc G}}
\]
where the composition in the \textit{definiens} (on the right) 
is composition in the category of tight zigzags.
Suppose $F$, $G$ and $H$ are contracted tight zigzags 
so that $F \circ G \circ H$ is defined.
To show that composition of contracted tight  zigzags 
is associative we need to show the following identity amongst contractions of composed zizgags in the category of tight zigzags.
\[
\contract{\contract{F \circ G} \circ H}
=\contract{F \circ \contract{G \circ H}}
\]
If we assume that $F$, $G$ and $H$ fit together snugly
so that $F \circ G \circ H$ in tight zigzags is
exactly  $F \cnc G \cnc H$. This makes life easier.

With this simplication we need to show that 
\[
\contract{\contract{F \cnc G} \cnc H}
=\contract{F \cnc \contract{G \cnc H}}.
\]
which is exactly lemma \lref{snugAssociativity}.

If on the otherhand $F$, $G$ and $H$ do not fit together snugly
then \highlight{...how to proceed}.